\documentclass[11pt,a4paper]{article}

\usepackage[utf8]{inputenc}
\usepackage[T1]{fontenc}
\usepackage{lmodern}
\usepackage[margin=2.5cm]{geometry}
\usepackage{microtype}
\usepackage{listings}
\usepackage{xcolor}
\usepackage{booktabs}
\usepackage{hyperref}
\usepackage{enumitem}
\usepackage{graphicx}

\hypersetup{
  colorlinks=true,
  linkcolor=blue!50!black,
  urlcolor=blue!50!black,
  citecolor=blue!50!black
}

% A simple style for vintage-machine code listings.
\lstdefinestyle{mic}{
  basicstyle=\ttfamily\footnotesize,
  numbers=none,
  frame=single,
  framerule=0.4pt,
  rulecolor=\color{black!40},
  backgroundcolor=\color{black!3},
  breaklines=true,
  showstringspaces=false,
  columns=fullflexible,
  keepspaces=true,
  commentstyle=\color{black!55}\itshape,
  literate={->}{$\rightarrow$}{2},
}

\lstset{style=mic}

\title{Recursive Towers of Hanoi on 1980s Educational Microcomputers \\
  \large A Tale of Two Architectures, Forty-Two Years Late}
\author{LambdaMikel \and Claude (Opus 4.7, Anthropic)}
\date{May 2026}

\begin{document}
\maketitle

\begin{abstract}
We present two implementations of the recursive Towers of Hanoi
algorithm targeting educational microcomputer systems from the early
1980s: the \emph{Kosmos CP1} (1983, Intel 8049 based) and the
\emph{Busch Microtronic 2090} (1981, TI~TMS1600 based). The CP1
implementation reads as a textbook example of recursion: the machine
provides indirect addressing, indirect jumps, and enough memory to
host a conventional call stack. The Microtronic implementation,
however, must contend with a platform that offers none of those
primitives---no stack, no indirect addressing, no indirect jumps, no
random-access memory beyond the register file, and a four-bit data
path. We show how recursion can nonetheless be realised on the
Microtronic by exploiting its rarely-used dual register bank as a
substrate for two hand-built stacks, and by replacing indirect
returns with a small marker-dispatch table. We compare the two
implementations as engineering artefacts and conclude with a brief
account of physically embodying the algorithm by having either host
drive an Arduino-based robotic arm via a simple four-bit protocol.
\end{abstract}

\section{Introduction}

The Towers of Hanoi is the textbook example of tree recursion. Its
implementation in any modern programming language is unremarkable;
the language's runtime stack handles the recursive descent and the
problem reduces to a six-line function. On a machine that has no
hardware stack, no indirect addressing, and only sixteen four-bit
registers, the story is rather different.

This article concerns two such machines, both sold as educational
microcomputers in the early 1980s, and both essentially obsolete by
1990:

\begin{itemize}
  \item the \textbf{Kosmos CP1} (1983), built around an Intel 8049
    microcontroller, with hexadecimal keypad and seven-segment display,
    and (with the CP3 expansion) 256 words of memory;
  \item the \textbf{Busch Microtronic 2090} (1981), built around a TI
    TMS1600 calculator-derived microcontroller, with the same physical
    interface but a fundamentally simpler architecture---no addressable
    memory other than the register file, and no stack of any kind.
\end{itemize}

Both machines were sold to teach young people how computers ``really
work,'' and both succeed at this in ways the marketing material did
not anticipate. The CP1 teaches you what \emph{a} computer looks like.
The Microtronic teaches you what a computer might be made to do when
the architecture refuses to help.

We will spend the bulk of this article on the Microtronic
implementation, because that is the more interesting story, and only
briefly summarise the CP1 program, which is in essence the form one
would write if given the choice of platform. We close with a short
discussion of the recent extension of both programs to physically
drive a Hanoi-solving robot through an Arduino-mediated digital
protocol, demonstrating that the original programs were really doing
useful work, not merely simulating it on a seven-segment display.

The recursive procedure under discussion is, throughout, the
following:

\begin{lstlisting}
function MoveTower(disk, source, dest, spare):
    if disk == 1:
        emit_move(source, dest)
    else:
        MoveTower(disk - 1, source, spare, dest)
        emit_move(source, dest)
        MoveTower(disk - 1, spare, dest, source)
\end{lstlisting}

\noindent
The challenge on either machine is implementing the two recursive
calls without losing the caller's local state. On a stack machine
this is free; on the CP1 it is a deliberate engineering exercise; on
the Microtronic it is essentially a small piece of architecture
research.

\section{The Kosmos CP1}

\subsection{What the machine provides}

The CP1 is a fairly orthodox 8-bit accumulator machine. With the CP3
memory expansion it offers 256 words of writable memory, addressable
0--255. Beyond the conventional load/store/arithmetic operations, it
offers three primitives that turn out to be decisive for any
recursive program:

\begin{itemize}
  \item \texttt{LIA} ($19.xxx$): \emph{load indirect}. Treats memory
    cell $xxx$ as a pointer; loads the accumulator from the address
    that cell contains.
  \item \texttt{AIS} ($20.xxx$): \emph{store indirect}. Stores the
    accumulator at the address pointed to by cell $xxx$.
  \item \texttt{SIU} ($21.xxx$): \emph{jump indirect}. Transfers
    control to the address held in cell $xxx$.
\end{itemize}

\noindent
With these three instructions, recursion is straightforward: one
implements a stack as a region of memory, manages a stack pointer in a
fixed cell, and uses \texttt{SIU} to return to the caller via a saved
address.

\subsection{Memory layout}

The CP1 implementation maintains four stacks simultaneously: three
value stacks (\textit{source}, \textit{spare}, \textit{dest}) and one
return stack. Rather than packing all four into a single
interleaved frame layout, the program uses three \emph{parallel} value
stacks at stride three, plus a separate return stack:

\begin{lstlisting}
  140  aux 1            150  const 1
  141  aux 2            151  const 3
  142  aux 3
  143  return slot 1
  144  return slot 2

  145  disk number
  146  *source          (stack pointer; initially 200)
  147  *spare           (stack pointer; initially 201)
  148  *dest            (stack pointer; initially 202)
  149  *return          (stack pointer; initially 240)

  200..239  value stacks  (13 frames, stride 3)
  240..255  return stack  (16 entries)
\end{lstlisting}

\noindent
The three value-stack pointers begin at 200, 201, and 202. Pushing a
frame is then a matter of adding three to each pointer, since the
three values for the new frame land naturally at the next stride-three
slot. The disk number is treated as a single mutable cell rather than
being stacked; it is decremented on entry and incremented on return,
which is correct because the algorithm only ever explores depths
monotonically.

\subsection{The subroutine-call convention}

The CP1 has no hardware call instruction. The program implements one
by hand: before any \texttt{SPU} (unconditional jump) used as a call,
the caller writes the desired return address into a fixed memory cell;
the callee returns by \texttt{SIU} through that cell. Two such cells
exist, named \emph{simple return 1} and \emph{simple return 2}, which
allows one level of subroutine nesting (the inner subroutine may
itself call another, provided it uses the other slot). For deeper
nesting---in particular, for recursion---the program pushes the return
address onto the return stack instead, and the dedicated
\textit{return block} at address 136 pops the top of the return stack
into \emph{simple return 1} and \texttt{SIU}s through it.

In effect, the CP1 program reconstructs the runtime apparatus
of a single-threaded recursive language: a call stack, a return
mechanism, and a frame layout. None of it is exotic. The interesting
choice is the \emph{stride-3 parallel stacks}, which save the
arithmetic of computing per-field offsets inside packed frames; once
the three pointers are advanced together, an indirect load of, say,
\texttt{*source} yields the source peg of the current frame without
further work.

\subsection{Code walkthrough}

The full recursive-call site, prologue, and epilogue are short enough
to read end-to-end. The first recursive call in \texttt{move\_tower}
performs \texttt{move\_tower(disk-1, source, spare, dest)}; its code
sits at cells 023--034:

\begin{lstlisting}
023 ako 04.026   # load 26 (= label_0)
024 abs 06.143   # write into 'simple return 1'
025 spu 09.098   # call save_and_push
026 lda 05.140   # accu = saved source
027 ais 20.146   # store into new top of source stack
028 lda 05.142   # accu = saved dest
029 ais 20.148   # store into new top of dest stack ...
                 # ... but used as the new SPARE: dest -> spare
030 lda 05.141   # accu = saved spare
031 ais 20.147   # store into new top of spare stack ...
                 # ... used as the new DEST: spare -> dest
032 ako 04.035   # load label_2 address
033 ais 20.149   # push onto return stack
034 spu 09.020   # GOTO move_tower
\end{lstlisting}

\noindent
Lines 028--031 are the swap that makes the recursion work: the new
\texttt{dest} is the old \texttt{spare}, and the new \texttt{spare}
is the old \texttt{dest}. The reads come from scratch cells
140/141/142 because \texttt{save\_and\_push} has already snapshotted
the caller's frame there before advancing the pointers.

The prologue, \texttt{save\_and\_push}, occupies cells 098--119:

\begin{lstlisting}
098 lda 05.145   # accu = n
099 sub 08.150   # n - 1
100 abs 06.145   # write back
101 lia 19.146   # accu = source (indirect load)
102 abs 06.140   #   ... save to scratch
103 lia 19.148   # accu = dest
104 abs 06.141   #   ... save to scratch
105 lia 19.147   # accu = spare
106 abs 06.142   #   ... save to scratch
107 lda 05.146   # advance *source by 3
108 add 07.151   #   (cell 151 holds 3)
109 abs 06.146
110 lda 05.148   # advance *dest by 3
111 add 07.151
112 abs 06.148
113 lda 05.147   # advance *spare by 3
114 add 07.151
115 abs 06.147
116 lda 05.149   # advance *return by 1
117 add 07.150   #   (cell 150 holds 1)
118 abs 06.149
119 siu 21.143   # return via 'simple return 1'
\end{lstlisting}

\noindent
The pattern is: snapshot the live frame into scratch, advance every
stack pointer, return. The call-site setup code, equipped with the
snapshot and the new (advanced) pointers, can now permute the
arguments and write them into the new top of each value stack
without needing temporaries of its own.

The epilogue, \texttt{pop\_and\_restore}, at cells 120--135 is the
mirror image: increment \texttt{n}, retreat all four stack pointers,
return. The frames themselves are left untouched in memory; the
restoration is purely a pointer reset, which is one of the small
elegances the indirect-addressing primitives \texttt{LIA}/\texttt{AIS}
buy on this machine.

Finally, the \emph{return block} at cells 136--138 does the indirect
jump that the Microtronic cannot:

\begin{lstlisting}
136 lia 19.149   # accu = top-of-return-stack (the saved label)
137 abs 06.143   # write into 'simple return 1'
138 siu 21.143   # indirect jump through 143 -- lands at the label
\end{lstlisting}

\noindent
Cell 143 is the same \emph{simple return 1} slot that
\texttt{save\_and\_push} also returns through; the return block
overwrites it with the popped address before jumping. The
\texttt{SIU} instruction is the irreducible primitive the
Microtronic lacks, and is what lets the CP1 implementation look,
overall, like a textbook recursive program.

\subsection{Capacity}

The CP1 program supports up to thirteen disks within the default
memory layout. With minor adjustments to stack base addresses (the
comments suggest moving the value stack start to 170 and the return
stack to 230), it can be pushed further still. The actual ceiling is
the available memory.

\section{The Busch Microtronic 2090}

\subsection{What the machine does not provide}

The Microtronic is a different proposition. Its architecture omits
nearly every primitive that makes the CP1 program straightforward:

\begin{itemize}
  \item There is no addressable memory beyond the register file. All
    state must live in registers.
  \item There are only sixteen registers, each four bits wide.
  \item There is no indirect addressing; all instructions encode their
    operands as register or constant nibbles.
  \item There is no indirect jump; control transfer is by absolute
    eight-bit address only.
  \item There \emph{is} a hardware \texttt{CALL}/\texttt{RET} pair,
    but it relies on a single return-address register that is neither
    user-accessible nor nestable. A second \texttt{CALL} issued while
    a return address is still live overwrites that register; the
    outer return address is then lost, and the subsequent
    \texttt{RET} jumps to the inner call site. The hardware
    \texttt{CALL}/\texttt{RET} mechanism is therefore confined to
    \emph{leaf} subroutines---it cannot, by itself, support
    recursion.
  \item There is, however, a curiosity that turns out to matter
    enormously: \emph{the register file is actually two banks of
    sixteen registers each}, with three exchange instructions
    (\texttt{EXRL}, \texttt{EXRM}, \texttt{EXRA}) that swap the
    contents of the low half, the high half, or the whole file with a
    hidden second bank.
\end{itemize}

\noindent
For most of the Microtronic's history, the second bank was used---if
at all---as a scratch area for temporarily preserving values across
some calculation. The implementation discussed here uses it for
something more structural.

\subsection{The central idea}

The implementation treats the Microtronic's two register banks as a
two-page memory system. One bank holds the current active frame and
the return stack; the other bank holds the value stack. The exchange
instructions \texttt{EXRL} and \texttt{EXRM}, applied together, swap
the two banks in their entirety; applied around a block of
register-shuffling code, they let one bank be mutated while the other
is left untouched and then ``swapped back into focus.''

Concretely:

\begin{itemize}
  \item \textbf{Bank A} (the bank active at procedure entry):
    registers~$\mathtt{R_8}$ to $\mathtt{R_F}$ hold the current frame
    and a return marker; registers $\mathtt{R_0}$ to $\mathtt{R_7}$
    are left alone.
  \item \textbf{Bank B} (accessed only via \texttt{EXRL}+\texttt{EXRM}
    pairs): all sixteen registers $\mathtt{R_0}$ to $\mathtt{R_F}$
    hold the value stack as a deque of up to four four-value frames.
\end{itemize}

\noindent
The shift routines that implement push and pop on the value stack
\emph{themselves} bracket their work with
\texttt{EXRL}+\texttt{EXRM}\,/\,\texttt{EXRL}+\texttt{EXRM}, so that
from the caller's perspective the shifts appear atomic and the
caller's own bank is restored on return. This bracketing pattern is
what makes the whole architecture composable, and it is what allows
the protocol code we describe in Section~\ref{sec:robot} to use the
otherwise-unused registers $\mathtt{R_0}$, $\mathtt{R_1}$,
$\mathtt{R_7}$ as scratch and persistent state without disturbing
either stack.

\subsection{Frame layout and the push}

Each call frame consists of four values---disk number, source peg,
destination peg, spare peg---occupying four consecutive registers.
Pushing a frame requires:

\begin{enumerate}
  \item Calling the \emph{value-stack shift-down} routine four times,
    each call shifting every register one position towards the lower
    indices (with the lowest register lost). After four shifts, four
    register slots have been freed at the top of the stack.
  \item Switching to bank~B (via \texttt{EXRM}) and writing the four
    new frame values into the freed slots.
  \item Switching back to bank~A, recording a small integer
    \emph{return marker} in $\mathtt{R_F}$ to identify the call site,
    and unconditionally jumping (\texttt{GOTO}) to the recursive entry
    point.
\end{enumerate}

\noindent
The shift routine itself is unrolled and minimal:

\begin{lstlisting}
70: f0d   # EXRL: swap low half of register file
71: f0e   # EXRM: swap high half
72: 010   # r0 := r1   (operand order: 0RS means r{S} := r{R})
73: 021   # r1 := r2
74: 032   # r2 := r3
75: 043   # r3 := r4
76: 054   # r4 := r5
77: 065   # r5 := r6
78: 076   # r6 := r7
79: 087   # r7 := r8
7a: 098   # r8 := r9
7b: 0a9   # r9 := rA
7c: 0ba   # rA := rB
7d: 0cb   # rB := rC
7e: 0dc   # rC := rD
7f: 0ed   # rD := rE
80: 0fe   # rE := rF   (rF is the only register not written; it duplicates into rE)
81: f0e   # EXRM (back)
82: f0d   # EXRL (back)
83: f07   # RET
\end{lstlisting}

\noindent
The two \texttt{EXR*} pairs are the magic. Inside the bracket the
machine is operating on bank~B, where the value stack lives; outside,
on bank~A, where the caller's working registers are. A single call to
this routine shifts the entire bank-B register file by one position
\emph{without} touching bank~A. Four calls in a row create the
four-slot hole that a push requires.

A complete recursive call site brings the pieces together. The
listing below is the program's first recursive call, performing
\texttt{MoveTower(n-1, source, spare, dest)}:

\begin{lstlisting}
15: f0e   # EXRM (to bank A, caller-visible side)
16: b70   # CALL value-stack shift-down (1 of 4)
17: b70   # ...
18: b70   # ...
19: b70   # 4 shifts: four-slot hole opens at top of bank B
1a: f0e   # EXRM (to bank B, where the new frame goes)
1b: 0bf   # MOV RB->RF: new n = caller's n
1c: 71f   # SUBI 1, RF: n := n - 1
1d: 0ae   # MOV RA->RE: new source = caller's source
1e: 09c   # MOV R9->RC: new spare  = caller's dest
1f: 08d   # MOV R8->RD: new dest   = caller's spare
20: f0e   # EXRM (back to bank A)
21: b50   # CALL return-stack shift-down: room for a new marker
22: 12f   # MOVI 2, RF: push return marker = 2
23: c0d   # GOTO MoveTower entry
\end{lstlisting}

\noindent
Two points are worth dwelling on. First, the four \texttt{CALL}s at
\$16--\$19 are leaf calls---they use the hardware return register,
never nest into one another, and so the single-level
\texttt{CALL}/\texttt{RET} mechanism noted earlier is no obstacle.
The recursive call itself, at \$23, is deliberately \emph{not} a
\texttt{CALL} but a plain \texttt{GOTO}; the would-be return
information is recorded in the software return stack at \$21--\$22,
not in the hardware return register, which is left undisturbed.

Second, the new frame's \texttt{source}, \texttt{dest}, and
\texttt{spare} fields are read from \texttt{R8}, \texttt{R9}, and
\texttt{RA} on lines \$1d--\$1f, not from the \texttt{RC}--\texttt{RE}
positions where they originally lived. The four shift-down calls have
relocated them four positions toward the lower indices, and the
setup code knows this---it reads them in their post-shift positions
and writes them into the new top frame, performing the
\textit{dest}\,$\leftrightarrow$\,\textit{spare} swap that the
recursive call requires. The second recursive call, at \$2e--\$3b,
follows the same shape with the opposite permutation and a different
return marker.

\subsection{The return mechanism}

The hardware \texttt{CALL}/\texttt{RET} pair cannot nest, and the
Microtronic has neither a user-accessible program counter nor an
indirect jump. So even if the program were willing to save the return
address by hand, there is nowhere to save it \emph{from} and no way
to jump back to it \emph{through}. After a recursive call returns,
control must somehow find its way back to one of three possible
continuation sites in the body of \texttt{MoveTower}: the point after
the first recursive call, the point after the second, or the
top-level caller in main. Conventional return-address pushing is
unworkable on all three counts.

The implementation solves this with a tiny \emph{return marker}: each
call pushes a small integer onto a one-register-deep return stack (in
fact several registers deep, allowing the call chain to nest as
deeply as the algorithm requires). After the call returns, control
transfers to a fixed \emph{dispatch block}:

\begin{lstlisting}
b0: 91f   # CMPI 1, RF        (marker == 1?)
b1: E0b   # BRZ to 0b         (return to main)
b2: 92f   # CMPI 2, RF        (marker == 2?)
b3: E24   # BRZ to 24         (return to inner site 2)
b4: 93f   # CMPI 3, RF        (marker == 3?)
b5: E3c   # BRZ to 3c         (return to inner site 3)
b6: F00   # HALT              (error: unknown marker)
\end{lstlisting}

\noindent
With only three distinct call sites in the program, three markers
suffice. The dispatch costs a handful of compare-and-branch
instructions per return, which is cheap relative to the
sixty-instruction cost of pushing and popping a value-stack frame.

At the chosen return target, the caller's frame is reinstated by
running the push sequence in reverse. After the dispatch at
\$b3 sends control to \$24:

\begin{lstlisting}
24: b60   # CALL return-stack shift-up: pop the consumed marker
25: b90   # CALL value-stack shift-up (1 of 4)
26: b90   # ...
27: b90   # ...
28: b90   # 4 shifts: caller's frame back at RC..RF in bank B
\end{lstlisting}

\noindent
After \$28, the caller's \texttt{n}, \texttt{source}, \texttt{dest}
and \texttt{spare} are once again the top frame of the value stack,
and execution resumes with the caller none the wiser.

\subsection{Capacity}

The Microtronic implementation supports up to four disks. The
constraint is the value stack: each frame is four registers wide, and
the deepest sensible level holds four frames simultaneously, exactly
filling bank~B's sixteen registers. A more aggressive packing---two
two-bit peg labels and a two-bit disk number into a single
register---would extend this to seven or eight disks, at the cost of
significant additional pack/unpack code.

\section{Comparison}

The two implementations solve the same problem on machines of broadly
similar physical scale (256 words of program memory, six-digit
displays, hex keypads, single-digit kilohertz clock speeds), but they
read as if they belonged to different decades.

\begin{center}
\small
\begin{tabular}{p{4.6cm}p{4.6cm}p{4.6cm}}
\toprule
 & \textbf{Kosmos CP1} & \textbf{Busch Microtronic 2090} \\
\midrule
CPU                       & Intel 8049 (8-bit)             & TI TMS1600 (4-bit interpreted) \\
Memory beyond regs        & 256 words (with CP3)            & none \\
Indirect load/store       & \texttt{LIA}, \texttt{AIS}      & not available \\
Indirect jump             & \texttt{SIU}                    & not available \\
Bitwise AND               & \texttt{UND}                    & not available \\
Number of registers       & 8-bit accumulator               & 16 four-bit registers, two banks \\
Stack realisation         & memory region + pointer cell    & shifts across the second bank \\
Subroutine return         & indirect via memory cell        & marker dispatch table \\
Disk-count ceiling        & 13 (default), more with effort  & 4 \\
Per-move overhead         & a few dozen instructions        & roughly sixty (mostly shifts) \\
Lines of source           & $\sim$140                       & $\sim$160 \\
Reads like                & a recursive program             & a small machine you built \\
\bottomrule
\end{tabular}
\end{center}

\bigskip

\noindent
The CP1 implementation is \emph{applied} computer science: it is what
one writes when the platform provides the primitives the algorithm
needs. The Microtronic implementation is \emph{experimental} computer
science: it is what one writes when those primitives must first be
invented out of the available substrate. The bank-swap-as-stack idiom
is the central piece of invention. It would generalise to other
recursive algorithms---tree walks, mutual recursion, even simple
interpreters---on the same hardware, but to the author's knowledge it
had not appeared in print, in forty-two years of the Microtronic's
existence, until this implementation was written in early 2024.

\section{Physical Embodiment}
\label{sec:robot}

Both implementations have been extended to drive a physical
Towers-of-Hanoi-solving robot. The robot is mechanically simple: a
two-degree-of-freedom pan/tilt arm with an electromagnet at the tip,
controlled by an Arduino. The Arduino performs all servo motion
planning, magnet pulse-width control, and disk-position tracking; the
vintage host need only emit a stream of high-level
\(\langle\textit{source},\textit{dest}\rangle\) move commands.

The communication protocol is deliberately minimal so that it can fit
inside either machine's tiny I/O envelope. Four output lines from the
host carry a three-bit move code (1--6, one per legal source-to-dest
pair) and a strobe bit that toggles on every command; one input line
returns to the host as a \emph{busy} signal that is held high while
the Arduino is in motion. The host's per-move logic is then:

\begin{enumerate}
  \item Poll the busy line until low.
  \item Compute the three-bit move code and toggle the strobe bit.
  \item Atomically present all four output bits.
  \item Poll the busy line until high (acknowledgement).
\end{enumerate}

\noindent
On the Microtronic, the four output bits are written by a single
\texttt{DOT} instruction, which makes the protocol naturally
race-free: the strobe transition and the data bits all change in the
same instruction. On the CP1, the same effect is achieved by writing
the three data bits first via three successive \texttt{P2A}
instructions and the strobe bit last, so that when the Arduino's
edge-triggered interrupt fires on the strobe, the data lines are
already settled.

A subtle question on the Microtronic was where to store the
single bit of persistent strobe state across recursive calls. The
algorithm uses every register in both banks for the value and return
stacks---or does it? In fact, registers $\mathtt{R_0}$ to
$\mathtt{R_7}$ of the active (bank~A) state are never written by the
original recursive code. The shift routines do touch them, but
\emph{only} inside their \texttt{EXRL}-bracketed regions, where they
operate on bank~B's copies; the active-bank values are restored on
exit. The strobe bit therefore lives in $\mathtt{R_7}$ of bank~A and
survives the full depth of the recursion without further machinery.

On the CP1, the equivalent state is simply a memory cell.

\section{Closing Remarks}

The CP1 program is a finished engineering artefact: it does what
recursion is supposed to do, on a machine that lets it. It is the
program one would write today if one were trying to teach a student
how recursion is implemented underneath a high-level language.

The Microtronic program is something different. It is a discovery
about a forty-year-old machine---a demonstration that an architecture
which appears to forbid recursion will, in fact, host it, provided
the programmer is willing to construct the missing primitives by
hand. The bank-swap-as-stack idea is small, but it is the kind of
thing that does not announce itself: in forty-two years of the
Microtronic's existence, it had not been written down. It is in the
nature of obscure platforms that they accumulate such secrets, and in
the nature of long acquaintance that one occasionally notices them.

That the same machines can now drive a physical robot, through a
fifty-cent serial protocol to a five-dollar microcontroller, is a
fitting late chapter. The TMS1600 and the 8049 were designed to teach
people what computers were; they are now, against considerable odds,
teaching one of them to move steel washers between three pegs in a
provably-optimal sequence. The teaching, it appears, was effective.

\section*{Authorship Note}

The prose of this paper was generated by Claude Opus~4.7 (Anthropic),
working from extended discussions of the source code with the
human co-author. The two recursive Towers of Hanoi
implementations themselves---the Microtronic 2090 program and the
Kosmos CP1 program---were written by hand by the human co-author
(LambdaMikel) and should be credited accordingly. The robot-control
extensions to both programs (the Arduino sketch, the patched
\texttt{HANOI-ROBOT.MIC} and \texttt{HANOIC-CP1-ROBOT.txt}, and the
shared four-bit GPIO protocol described in
Section~\ref{sec:robot}) were designed and written by Claude during
the conversation that produced this paper.

\appendix

\section{Source-Code Availability}

The reference implementations for both machines, and the Arduino
firmware for the robot, are available at
\url{https://github.com/lambdamikel/towers-of-hanoi}. The opcode
reference used in decoding the Microtronic source is drawn from the
emulator code at \url{https://github.com/lambdamikel/Busch-2090}; the
Kosmos CP1 manual is available via the Internet Archive at
\url{https://archive.org/details/cp-1-manual}.

\end{document}
